{\color{blue}\paragraph{Positioning of this study}
Closest to our work are graph-based DRL portfolio models such as GPM~\citep{Si2022GPM}, multi-GCGRU~\citep{Liu2021MultiGCGRU}, and DeepPocket~\citep{Farzan2021DeepPocket}, which embed relational information about stocks into the portfolio decision model, and graph-attention forecasters such as HATS~\citep{Kim2019HATS}, STGAT~\citep{Feng2025STGAT}, and BiLSTM-GAT~\citep{Lu2025BiLSTMGAT}, which weight related stocks by learned attention. MGAT-FinRL differs in three respects. First, its relation layers are built from qualitative financial ties (common top-ten shareholders and co-mentions in investor-relations questions) in addition to a trading-activity correlation, and they restrict, rather than weight, the attention between stocks; in the present data, the trading-value layer determines most of the admissible pairs (Section~\ref{subsec3}). Second, the forecaster is trained for the return that a decision can actually earn under next-day execution, and its forecasts enter the policy's state at every decision; forecasting and policy learning remain separate stages, so the forecast source can be varied while the policy optimiser, environment, and seeds are held fixed. Third, the evaluation is pre-registered and uses a year that entered no earlier analysis, with a naive forecast source of identical state dimension as the reference. The individual GAT, Transformer, and actor-critic components are standard; the contribution lies in the relation construction, the tradable-horizon objective, and a controlled evaluation of how much relational forecasts change portfolio outcomes.\par}
